% Local macros for cross-volume conductor references.
\providecommand{\BRST}{\mathrm{BRST}}
\providecommand{\BRSTGauged}{\mathrm{BRSTGaugedChirAlg}}
\providecommand{\PhiFA}{\Phi^{\mathrm{FA}}}
\providecommand{\SpCh}{\mathrm{Sp}^{\mathrm{ch}}}
\providecommand{\BP}{\mathrm{BP}}
\providecommand{\ClaimStatusComputed}{\textnormal{[computed]}}

\section{Verified Drinfeld--Sokolov and \texorpdfstring{$\mathcal W$}{W}-algebra lanes}
\label{sec:walg-deep-repaired}

The universal affine \(\mathcal W\)-algebra
\(\mathcal W^k(\fg,f)\) is the degree-zero cohomology of the quantum
Drinfeld--Sokolov BRST complex attached to a good grading for the
nilpotent element \(f\).  This chapter uses that construction at
Beilinson level~\(1\), its ordered chiral bar complex at level~\(2\),
and scalar genus-one data at level~\(5\).  Each passage between these
levels carries its own comparison map.

Four objects occur throughout:
\[
\mathcal W^k(\fg,f),\qquad
\bar B_X\!\left(\mathcal W^k(\fg,f)\right),\qquad
\left(\mathcal W^k(\fg,f)\right)^{\mathrm i},\qquad
\left(\mathcal W^k(\fg,f)\right)^!.
\]
They denote respectively the chiral algebra, its chain-level ordered
bar coalgebra, its bar cohomology, and a Verdier--Koszul dual algebra
after finite-type or completed duality.  A theorem comparing any two
members of this list states the required filtration, convergence, and
duality hypotheses.

\begin{remark}[Three verified surfaces]
\label{rem:w-deep-three-pillar}
\index{W-algebras@$\mathcal W$-algebras!verified surfaces}
The chapter has three mathematical surfaces.
\begin{enumerate}[label=\textup{(\roman*)}]
\item Quantum Hamiltonian reduction supplies the BRST complex and the
  universal algebra \(\mathcal W^k(\fg,f)\).
\item OPE calculations supply exact rational functions, including the
  principal central-charge companion and the BP companion sum.
\item Bar, Verdier, modular, and large-rank comparisons are derived
  statements whose hypotheses are recorded explicitly.
\end{enumerate}
The principal \(W_N\) scalar formulas and the Bershadsky--Polyakov
central formula therefore occupy separate theorem lanes.
\end{remark}

\subsection{The derived Drinfeld--Sokolov functor}
\label{sec:ds-derived-foundation}

Fix a good \(\frac12\mathbb Z\)-grading
\(\fg=\bigoplus_j\fg_j\) for \(f\), with charged ghosts for
\(\fg_{>0}\) and neutral fields for \(\fg_{1/2}\).  We write
\(F^{\mathrm{ch}}_f\) and \(F^{\mathrm{ne}}_f\) for the corresponding
free-field factors.

\begin{definition}[Filtered derived Drinfeld--Sokolov functor;
\ClaimStatusDefinitional]
\label{def:filtered-derived-ds-functor}
\index{Drinfeld--Sokolov reduction!derived functor|textbf}
Let \(A\) be an augmented chiral algebra with a compatible
\(\widehat\fg_k\)-action and a complete, separated Kazhdan
filtration, bounded below in every conformal-weight window.  Define
\[
\mathbf R\mathrm{DS}_f(A)
:=
\operatorname{Tot}\!\left(
A\widehat\otimes F^{\mathrm{ch}}_f
\widehat\otimes F^{\mathrm{ne}}_f,\,
Q_{\mathrm{DS}}
\right).
\]
The differential is the Kac--Roan--Wakimoto quantum reduction
differential, including the current, charged-ghost, neutral-field,
and character terms \cite{KRW}.  A morphism \(u\colon A\to A'\)
induces
\(\operatorname{Tot}(u\widehat\otimes\mathrm{id})\).
The ordinary reduction is
\[
\mathrm{DS}_f(A):=H^0\!\left(\mathbf R\mathrm{DS}_f(A)\right).
\]
The equality of derived and degree-zero reduction is governed by the
BRST concentration condition
\[
H^i\!\left(\mathbf R\mathrm{DS}_f(A)\right)=0
\quad (i\ne0).
\]
\end{definition}

\begin{definition}[Screening operators]
\label{def:screening}
\ClaimStatusDefinitional
\index{screening operator|textbf}
In a chosen free-field realization, a screening current \(S_i(z)\)
is a field whose residue
\[
Q_i:=\operatorname{Res}_{z=0}S_i(z)
\]
commutes with the \(\mathcal W\)-action up to a total derivative.
An intersection-of-kernels presentation
\(\mathcal W=\bigcap_i\ker Q_i\) is a theorem about that realization,
level, and completion.  A screening resolution additionally requires
exactness in every nonzero cohomological degree.
\end{definition}

\begin{remark}[BRST concentration and screening exactness]
Kac--Roan--Wakimoto construct the quantum reduction complex
\cite{KRW}.  Free-field screening presentations supply an effective
description on generic Wakimoto lanes.  Critical, admissible, and
simple-quotient fibers carry their own concentration and exactness
statements.
\end{remark}

\subsection{Principal type \(A\): the exact scalar lane}
\label{sec:principal-wn-exact}

For \(\fg=\mathfrak{sl}_N\), let \(f_{\mathrm{prin}}\) be principal
and write \(\mathcal W_N^k=\mathcal W^k(\mathfrak{sl}_N,f_{\mathrm{prin}})\).
The universal principal algebra at generic level has strong
generators of weights \(2,3,\ldots,N\).

\begin{theorem}[Fateev--Lukyanov central charge and reflected sum]
\label{thm:wn-central-companion-deep}
\ClaimStatusProvedHere
\emph{Type signature:}
\((\mathrm{Open},\mathrm{universal\ principal\ DS},1,
H_{\mathrm{univ}}^{W_N})\).
For \(N\ge2\) and \(k\ne-N\), the principal central charge is
\begin{equation}
\label{eq:wn-central-deep}
c_N(k)
=(N-1)-N(N^2-1)\frac{(k+N-1)^2}{k+N}.
\end{equation}
Under the reflected parameter \(k^\vee=-k-2N\),
\begin{equation}
\label{eq:wn-central-conductor-deep}
c_N(k)+c_N(k^\vee)
=4N^3-2N-2.
\end{equation}
Equation~\eqref{eq:wn-central-deep} is the Fateev--Lukyanov
principal-DS formula \cite[Equation~\textup{(}7\textup{)}]{FateevLukyanov88}.
\end{theorem}

\begin{proof}
Put \(t=k+N\) and \(A_N=N(N^2-1)\).  Then
\[
c_N(t)=(N-1)-A_N\left(t-2+t^{-1}\right).
\]
The reflection sends \(t\) to \(-t\), so
\[
c_N(t)+c_N(-t)
=2(N-1)+4A_N
=4N^3-2N-2.
\]
\end{proof}

\begin{computation}[The reflection in the shifted parameter]
\label{comp:ff-thooft-involution}
\ClaimStatusComputed
The parameter \(t=k+N\) linearizes the reflected companion:
\[
t^\vee=k^\vee+N=-t.
\]
This identity proves the scalar sum
\eqref{eq:wn-central-conductor-deep}.  Its promotion to an algebra
equivalence requires a bar--Verdier comparison theorem.
\end{computation}

\begin{proposition}[Principal genus-one scalar package]
\label{prop:higher-w-kappa-matrix}
\ClaimStatusConditional
\emph{Type signature:}
\((\mathrm{Open},\mathrm{principal\ DS},5,
H_{\mathrm{genus1}}^{W_N}\cup H_{\mathrm{comp}}^{W_N})\).
Assume the principal genus-one curvature theorem
\(H_{\mathrm{genus1}}^{W_N}\), normalized so that a weight-\(s\)
principal generator contributes \(c/s\), and the reflected companion
package \(H_{\mathrm{comp}}^{W_N}\).  Then
\[
\kappa(\mathcal W_N^k)
=c_N(k)\sum_{s=2}^{N}\frac1s
=c_N(k)(H_N-1),
\qquad
H_N:=\sum_{j=1}^{N}\frac1j,
\]
and the pointed Theorem~D class satisfies
\[
 \operatorname{tr}_1
 \operatorname{Obs}^{\mathrm{def}}_{1,1}(\mathcal W_N^k)
 =\kappa(\mathcal W_N^k)\lambda_1.
\]
and
\begin{equation}
\label{eq:wn-kappa-conductor-deep}
K_N^\kappa
=(H_N-1)(4N^3-2N-2).
\end{equation}
The diagonal OPE-normalization matrix
\[
D_N(c):=\operatorname{diag}\left(\frac c2,\frac c3,\ldots,\frac cN\right)
\]
has trace \(c(H_N-1)\).  The theorem identifies this trace with the
genus-one scalar precisely through
\(H_{\mathrm{genus1}}^{W_N}\).
\end{proposition}

\begin{proof}
The first equality is the stated genus-one hypothesis.  Add the two
reflected values and apply
\eqref{eq:wn-central-conductor-deep}.
\end{proof}

\begin{center}
\renewcommand{\arraystretch}{1.2}
\begin{tabular}{cccc}
\toprule
\(N\) & \(H_N-1\) & \(K_N^c\) & \(K_N^\kappa\) \\
\midrule
2 & \(1/2\) & \(26\) & \(13\) \\
3 & \(5/6\) & \(100\) & \(250/3\) \\
4 & \(13/12\) & \(246\) & \(533/2\) \\
5 & \(77/60\) & \(488\) & \(9394/15\) \\
\bottomrule
\end{tabular}
\end{center}

\begin{computation}[The \(\mathfrak{sl}_3\) DS hierarchy]
\label{comp:sl3-ds-hierarchy}
\ClaimStatusComputed
The three nilpotent partitions are
\[
(1,1,1),\qquad(2,1),\qquad(3).
\]
They index respectively the affine algebra, the
Bershadsky--Polyakov reduction, and the principal
\(\mathcal W_3\) reduction.  The partition labels organize the
reductions.  Verdier--Koszul companions require the comparison
package of \S\ref{sec:ds-bar-comparison-deep}.
\end{computation}

\begin{computation}[The \(\mathfrak{sl}_4\) orbit list]
\label{comp:sl4-ds-hierarchy}
\ClaimStatusComputed
The nilpotent partitions are
\[
(1^4),\quad(2,1^2),\quad(2^2),\quad(3,1),\quad(4).
\]
Their transposes give a combinatorial involution on the orbit list.
A functorial duality between the associated affine
\(\mathcal W\)-algebras is a separate theorem-level obligation.
\end{computation}

\subsection{Miura and \(W_3\) OPE calculations}
\label{sec:w3-verified-ope}

\begin{computation}[Quantum Miura operator]
\label{comp:miura-w3}
\ClaimStatusProvedElsewhere
Let \(\varphi_1+\cdots+\varphi_N=0\) be Heisenberg fields in the
standard principal free-field realization.  The normally ordered
operator
\[
:\prod_{i=1}^{N}
\left(\partial+\alpha_0\,\partial\varphi_i\right):
=\partial^N+\sum_{s=2}^{N}U^{(s)}\partial^{N-s}
\]
produces the Miura generators.  A triangular correction by
derivatives and normally ordered lower-spin fields yields primary
generators \(W^{(s)}\).  For \(N=3\), this gives the stress tensor
\(T\) and a weight-three primary \(W\).
\end{computation}

\begin{computation}[Universal \(W_3\) OPE packet]
\label{comp:w-infty-shadow-tower}
\label{comp:w3-shadow-tower-data}
\ClaimStatusProvedElsewhere
In Zamolodchikov normalization \cite{Zamolodchikov1985}, set
\[
\Lambda={:}TT{:}-\frac{3}{10}\partial^2T.
\]
The singular part of the universal \(W_3\) OPE is
\begin{align}
W(z)W(w)\sim{}&
\frac{c/3}{(z-w)^6}
+\frac{2T(w)}{(z-w)^4}
+\frac{\partial T(w)}{(z-w)^3}
\nonumber\\
&+\frac{\frac{32}{5c+22}\Lambda(w)
+\frac{3}{10}\partial^2T(w)}{(z-w)^2}
+\frac{\frac{16}{5c+22}\partial\Lambda(w)
+\frac1{15}\partial^3T(w)}{z-w}.
\label{eq:w3-wline-closed-form}
\end{align}
This finite OPE packet supplies exact collision coefficients.
Translation covariance fixes the pole-one coefficient to one half of
the pole-two \(\Lambda\)-coefficient.  If the latter is \(A\), these
two poles contribute \(\frac{A}{2}(m-n)\Lambda_{m+n}\) to the mode
commutator.  Zamolodchikov's coefficient
\(16(m-n)/(22+5c)\) therefore gives \(A=32/(22+5c)\), providing an
independent normalization check.
An all-degree shadow recurrence requires the transferred modular
deformation complex.
\end{computation}

\begin{proposition}[Weight-four Virasoro Gram matrix]
\label{prop:gram-wt4}
\ClaimStatusProvedHere
In the vacuum Verma module with basis
\[
L_{-4}\lvert0\rangle,\qquad L_{-2}^2\lvert0\rangle,
\]
the Gram matrix and its determinant are
\[
G_4=
\begin{pmatrix}
5c&3c\\[2pt]
3c&c(c+8)/2
\end{pmatrix},
\qquad
\det G_4=\frac{c^2(5c+22)}2.
\]
\end{proposition}

\begin{proof}
The Virasoro relation gives
\[
\langle0|L_4L_{-4}|0\rangle=5c,\qquad
\langle0|L_4L_{-2}^2|0\rangle=3c,
\]
and
\[
L_2L_{-2}^2|0\rangle=(c+8)L_{-2}|0\rangle,
\qquad
\langle0|L_2^2L_{-2}^2|0\rangle=\frac{c(c+8)}2.
\]
Taking the determinant gives the stated factor.
\end{proof}

\begin{corollary}[The weight-four quasi-primary]
\label{cor:lambda-qp}
\ClaimStatusProvedHere
The state
\[
\lvert\Lambda\rangle
=L_{-2}^2\lvert0\rangle-\frac35L_{-4}\lvert0\rangle
\]
satisfies \(L_1\lvert\Lambda\rangle=0\) and
\[
\langle\Lambda|\Lambda\rangle=\frac{c(5c+22)}{10}.
\]
Under state--field correspondence it is
\(\Lambda={:}TT{:}-\frac3{10}\partial^2T\).
\end{corollary}

\begin{theorem}[\(\mathcal W_4\) primary-normalization obligation]
\label{thm:c334}
\ClaimStatusOpen
Fix a primary generator \(W^{(4)}\), its two-point normalization,
and a universal \(\mathcal W_4\) presentation.  The coefficient of
\(W^{(4)}(w)(z-w)^{-2}\) in
\(W^{(3)}(z)W^{(3)}(w)\) is then a rational function of \(c\).
A proved formula requires a primary-source Miura or Jacobi
calculation with these normalizations and a quotient-by-quotient
treatment at determinant zeros.  The weight-four Gram factor above
supplies one input to that calculation.
\end{theorem}

\subsection{The Bershadsky--Polyakov truth surface}
\label{sec:bp-truth-surface}

Write
\[
\mathcal B^k
:=\mathcal W^k(\mathfrak{sl}_3,f_{(2,1)}),
\qquad k\ne-3.
\]
Fehily--Kawasetsu--Ridout define \(\mathcal B^k\) as an ordinary
bosonic vertex algebra \cite[Definition~2.1]{FKR20}.  Its strong
generators are
\[
J_{1}^{\mathrm{even}},\qquad
G^+_{3/2}^{\mathrm{even}},\qquad
G^-_{3/2}^{\mathrm{even}},\qquad
T_{2}^{\mathrm{even}}.
\]

\begin{proposition}[BP OPE normalization]
\label{prop:bp-ope-deep}
\ClaimStatusProvedElsewhere
In the Feigin--Semikhatov normalization used by
Fehily--Kawasetsu--Ridout,
\begin{align*}
J_{(1)}J&=\frac{2k+3}{3},&
J_{(0)}G^\pm&=\pm G^\pm,\\
T_{(1)}G^\pm&=\frac32G^\pm,&
G^+_{(2)}G^-&=(k+1)(2k+3),\\
G^+_{(1)}G^-&=3(k+1)J,&
G^+_{(0)}G^-&=
3{:}JJ{:}+\frac32(k+1)\partial J-(k+3)T.
\end{align*}
The Virasoro OPE has central charge
\begin{equation}
\label{eq:bp-central-deep}
c_{\mathrm{BP}}(k)
=-\frac{(2k+3)(3k+1)}{k+3}.
\end{equation}
\end{proposition}

\begin{theorem}[BP central companion]
\label{thm:bp-central-companion-deep}
\ClaimStatusProvedHere
\emph{Type signature:}
\((\mathrm{Open},\mathrm{FKR\ bosonic\ BP},5,
H_{\mathrm{central}}^{\mathrm{BP}})\).
For the reflected parameter \(k^\vee=-k-6\),
\begin{equation}
\label{eq:bp-central-sum-deep}
c_{\mathrm{BP}}(k)+c_{\mathrm{BP}}(k^\vee)=50
\qquad\text{in }\mathbb Q(k).
\end{equation}
\end{theorem}

\begin{proof}
Put \(t=k+3\).  Then
\[
c_{\mathrm{BP}}(k)=25-6t-\frac{24}{t},
\qquad
c_{\mathrm{BP}}(-k-6)=25+6t+\frac{24}{t}.
\]
Their sum is \(50\).
\end{proof}

\begin{proposition}[Secondary shifted BP scalar]
\label{prop:bp-shifted-secondary-deep}
\ClaimStatusComputed
The rational function
\[
c_{\mathrm{BP}}^{\mathrm{shift}}(k)
:=2-\frac{24(k+1)^2}{k+3}
\]
belongs to a shifted conformal-vector convention.  It satisfies
\[
c_{\mathrm{BP}}^{\mathrm{shift}}(k)
+c_{\mathrm{BP}}^{\mathrm{shift}}(-k-6)=196.
\]
Thus \(196\) records the secondary shifted lane, while
\eqref{eq:bp-central-sum-deep} records the canonical FKR lane.
\end{proposition}

\begin{proof}
With \(t=k+3\),
\[
c_{\mathrm{BP}}^{\mathrm{shift}}(k)
=98-24t-\frac{96}{t};
\]
the reflection \(t\mapsto-t\) gives the sum \(196\).
\end{proof}

\begin{proposition}[BP reciprocal-weight diagnostic]
\label{prop:bp-reciprocal-diagnostic-deep}
\ClaimStatusComputed
The even strong-generator packet gives
\[
1+\frac{1}{3/2}+\frac{1}{3/2}+\frac12
=1+\frac23+\frac23+\frac12
=\frac{17}{6}.
\]
This number is a reciprocal-weight diagnostic.  The BP genus-one
curvature is a separate BRST/modular computation.
\end{proposition}

\begin{computation}[BP bar and modular ledger]
\label{comp:bp-bar}
\ClaimStatusOpen
The presently computed BP data are
\[
\left(
\text{parity},\ \text{strong weights},\ \text{OPE packet},\
c_{\mathrm{BP}}(k),\ K^c_{\mathrm{BP}}
\right)
=
\left(
\mathrm{even},\ \{1,\tfrac32,\tfrac32,2\},\
\text{\eqref{eq:bp-central-deep} and the displayed OPEs},\
c_{\mathrm{BP}}(k),\ 50
\right).
\]
The following entries remain theorem obligations:
\[
\kappa_{\mathrm{BP}}(k),\qquad
\rho_{\mathrm{BP}}(k),\qquad
K^\kappa_{\mathrm{BP}},\qquad
(\mathcal B^k)^!,\qquad
\mathbf R\mathrm{DS}_{(2,1)}\bar B_X
\Longrightarrow
\bar B_X\mathbf R\mathrm{DS}_{(2,1)}.
\]
Their evaluation requires the genus-one BP BRST curvature and the
DS--bar comparison package below.
\end{computation}

\begin{remark}[Orbit label and duality datum]
\label{rem:bv-orbit-identification}
\index{Bershadsky--Polyakov algebra!orbit label}
The partition \((2,1)\) is self-transpose.  This supplies a
same-orbit candidate for a Verdier companion.  A same-family algebra
equivalence additionally requires a bar comparison, continuous
Verdier duality, cobar convergence, and a level-identification
theorem.
\end{remark}

\begin{proposition}[Partition-dependent complementarity obligation]
\label{prop:partition-dependent-complementarity}
\ClaimStatusOpen
For principal type \(A\), Theorem~\ref{thm:wn-central-companion-deep}
computes the reflected central sum.  For the BP partition,
Theorem~\ref{thm:bp-central-companion-deep} computes the value \(50\).
A formula for arbitrary partitions \(\lambda\vdash N\) requires the
KRW central charge in a fixed grading together with a proved
companion operation on nilpotent orbits.  The modular sum further
requires the genus-one curvature of each reduced family.
\end{proposition}

\subsection{DS, bar, and Verdier comparison}
\label{sec:ds-bar-comparison-deep}
\label{sec:ds-bar-intertwining}

\begin{proposition}[Filtered criterion for DS--bar exchange]
\label{prop:ds-bar-formality}
\ClaimStatusConditional
Let \(A\) satisfy the filtered-complete hypotheses of
Definition~\ref{def:filtered-derived-ds-functor}.  Assume:
\begin{enumerate}[label=\textup{(D\arabic*)}]
\item \(\mathbf R\mathrm{DS}_f\) has the lax monoidal structure needed
  to construct a filtered comparison
  \[
  \eta_{A,f}\colon
  \mathbf R\mathrm{DS}_f\!\left(\bar B_X(A)\right)
  \longrightarrow
  \bar B_X\!\left(\mathbf R\mathrm{DS}_f(A)\right);
  \]
\item both filtrations are complete, separated, and bounded below in
  every conformal-weight window;
\item \(\operatorname{gr}\eta_{A,f}\) is a quasi-isomorphism;
\item the relevant spectral sequences converge strongly.
\end{enumerate}
Then \(\eta_{A,f}\) is a filtered quasi-isomorphism.
\end{proposition}

\begin{proof}
Hypothesis~\textup{(D3)} identifies the \(E_1\)-pages.
Hypotheses~\textup{(D2)} and~\textup{(D4)} identify the abutments by
the comparison theorem for complete filtered complexes.
\end{proof}

\begin{remark}[Visible comparison obligations]
\label{rem:ds-bar-verdier-obstructions}
The verification package consists of the monoidal BRST comparison,
higher-cohomology control on every bar term, completeness on
Fulton--MacPherson strata, and Verdier dualizability.  Critical and
simple-quotient fibers change this package
through centers and null fields.  Non-principal reductions add the
neutral and ghost--ghost contributions to the associated graded.
\end{remark}

\begin{proposition}[DS-transferred Verdier--Koszul dual]
\label{prop:ds-transferred-koszul-dual}
\ClaimStatusConditional
Assume Proposition~\ref{prop:ds-bar-formality}, BRST concentration,
and Verdier duality through BRST totalization.  Then the
DS-transferred Verdier--Koszul algebra is
\[
\left(\mathbf R\mathrm{DS}_f(A)\right)^!_{\mathrm{DS}}
:=
\mathbb D_{\Ran}
\bar B_X\!\left(\mathbf R\mathrm{DS}_f(A)\right),
\]
and the comparison gives
\[
\mathbb D_{\Ran}\bar B_X\!\left(\mathbf R\mathrm{DS}_f(A)\right)
\simeq
\mathbf R\mathrm{DS}_f^\vee
\!\left(\mathbb D_{\Ran}\bar B_X(A)\right).
\]
A strict same-family \(\mathcal W\)-algebra presentation is obtained
by identifying the right-hand Verdier algebra through an independent
presentation theorem.
\end{proposition}

\begin{theorem}[DS/bar master square]
\label{thm:master-commutative-square}
\ClaimStatusConditional
Under hypotheses \textup{(D1)--(D4)} of
Proposition~\ref{prop:ds-bar-formality}, the square
\[
\begin{tikzcd}[column sep=large,row sep=large]
\mathbf R\mathrm{DS}_f(\bar B_X(A))
  \arrow[r,"\eta_{A,f}"] \arrow[d,"\mathbb D_{\Ran}"']
&
\bar B_X(\mathbf R\mathrm{DS}_f(A))
  \arrow[d,"\mathbb D_{\Ran}"]
\\
\mathbb D_{\Ran}\mathbf R\mathrm{DS}_f(\bar B_X(A))
  \arrow[r]
&
\mathbb D_{\Ran}\bar B_X(\mathbf R\mathrm{DS}_f(A))
\end{tikzcd}
\]
commutes in the filtered derived category once the lower horizontal
map is supplied by Verdier compatibility.  This is the comparison
surface used by any DS-transferred Koszul duality statement.
\end{theorem}

\begin{computation}[The \(\mathfrak{sl}_3\to W_3\) comparison window]
\label{comp:ds-bar-sl3-w3}
\ClaimStatusOpen
The first proof window must construct \(\eta_{A,f_{\mathrm{prin}}}\)
on bar lengths \(1,2,3\), retain the charged and neutral BRST fields,
and verify compatibility with every collision stratum.  The affine
Jacobi relation supplies one associated-graded cancellation.  The
remaining BRST/bar terms form the finite chain-map calculation.
\end{computation}

\begin{conjecture}[DS--Verdier intertwining for arbitrary nilpotent]
\label{conj:ds-kd-arbitrary-nilpotent}
\ClaimStatusConjectured
For a good grading of \(f\), assume windowwise perfectness,
BRST concentration, strong convergence of the Kazhdan and bar
spectral sequences, and Verdier compatibility on every
Fulton--MacPherson stratum.  Then
\[
\left(\mathbf R\mathrm{DS}_f(A)\right)^!_{\mathrm{DS}}
\simeq
\mathbf R\mathrm{DS}_f^\vee(A^!)
\]
in the completed derived category.
\end{conjecture}

\begin{remark}[Springer and Barbasch--Vogan input]
\label{rem:bv-springer-duality}
Springer and Barbasch--Vogan correspondences organize orbit data.
Their use in chiral Koszul duality requires a functor from the
geometric correspondence to the BRST/bar category and a theorem
identifying its output with the completed cobar algebra.
\end{remark}

\begin{remark}[Abelian and non-abelian nilradicals]
\label{rem:abelian-nonabelian-nilradical}
An abelian positive nilradical removes the ghost--ghost bracket term
from \(Q_{\mathrm{DS}}\).  A non-abelian nilradical retains it.
This distinction controls one part of the filtered differential;
orbit companionship and Verdier duality are additional structures.
\end{remark}

\begin{theorem}[Type-\(A\) transport-closure obligation]
\label{thm:transport-closure-type-a}
\ClaimStatusOpen
A transport-closure theorem for all partitions of \(N\) consists of:
the orbit correspondence, the level map, the filtered DS--bar
quasi-isomorphism, Verdier compatibility, and a presentation theorem
for the resulting cobar algebra.  Principal and hook calculations
provide test families for this package.  The arbitrary-partition
closure remains an open theorem lane.
\end{theorem}

\subsection{Degree-three ordered bar calculation for \(W_3\)}
\label{sec:w3-bar-degree3}

Let \(\bar B_{X,\le3}(\mathcal W_3)\) denote the ordered chiral bar
complex truncated at bar length three.  Its differential combines
the full singular OPE with the boundary maps of compactified
configuration spaces.  Mode-zero Chevalley--Eilenberg data form one
associated-graded piece.

\begin{computation}[Length-three chain groups]
\label{comp:w3-deg3-structure}
\ClaimStatusComputed
The length-three summand has the form
\[
\bigl(s^{-1}\overline{\mathcal W}_3\bigr)^{\boxtimes3}
\otimes
C^\bullet_{\mathrm{dR}}\!\left(\operatorname{Conf}_3(X)\right),
\]
with coinvariants and completions fixed by the ordered-bar
construction.  Conformal-weight truncation makes this a finite
linear-algebra problem in each window.
\end{computation}

\begin{computation}[Mixed \(T,W\) sectors]
\label{comp:w3-deg3-mixed}
\ClaimStatusComputed
The \(T/W\) colorings at length three are
\[
TTT,\quad TTW,\quad TWT,\quad WTT,\quad
TWW,\quad WTW,\quad WWT,\quad WWW.
\]
Their differentials are determined by
\eqref{eq:w3-wline-closed-form}, the Virasoro OPE, and the
configuration-space boundary signs.
\end{computation}

\begin{computation}[Arnold relation]
\label{comp:w3-arnold-deg3}
\ClaimStatusProvedElsewhere
On \(\operatorname{Conf}_3(X)\), the logarithmic one-forms satisfy
\[
\omega_{12}\wedge\omega_{23}
+\omega_{23}\wedge\omega_{31}
+\omega_{31}\wedge\omega_{12}=0.
\]
This identity supplies the geometric cancellation in the
length-three boundary differential.
\end{computation}

\begin{proposition}[Degree-three vacuum cohomology obligation]
\label{prop:w3-deg3-vacuum}
\ClaimStatusOpen
The value of the vacuum-sector cohomology in a conformal-weight
window is the kernel modulo image of the complete matrix assembled
from the OPE and Fulton--MacPherson boundary maps.  Row reduction of
that matrix is the required computation.
\end{proposition}

\begin{computation}[Degree-three cohomology ledger]
\label{comp:w3-deg3-cohom}
\ClaimStatusOpen
The ledger records chain bases, differential matrices, ranks,
kernel generators, and boundary representatives.  The present
chapter supplies the exact OPE input and the Arnold relation.
The rank computation remains the open entry.
\end{computation}

\begin{computation}[Reflected \(W_3\) scalar]
\label{comp:w3-curvature-dual}
\ClaimStatusConditional
Under the principal genus-one package,
\[
\kappa(\mathcal W_3^k)=\frac56c_3(k),\qquad
\kappa(\mathcal W_3^{k^\vee})=\frac56(100-c_3(k)),
\]
so the reflected scalar sum is \(250/3\).  The algebra-level
companion interpretation uses the DS--bar--Verdier package.
\end{computation}

\subsection{Modular Koszul packages}

\begin{definition}[Modular Koszul triple]
\label{def:modular-koszul-triple}
\ClaimStatusDefinitional
A modular Koszul triple is
\[
\mathfrak T_A=
\left(A,\,(A^{\mathrm i},A^!),\,r_A^{\mathrm{coll}}\right)
\]
together with:
\begin{enumerate}[label=\textup{(\roman*)}]
\item an augmented quadratic chiral algebra
  \(A=T_X(V)/(R)\);
\item its presentation coalgebra
  \(A^{\mathrm i}=C_X(s^{-1}V,s^{-2}R)\), full bar chain model
  \(\bar B_X(A)\), and comparison
  \(q_A\colon A^{\mathrm i}\to\bar B_X(A)\);
\item the Verdier algebra
  \(K_X(A)=\mathbb D_{\Ran}\bar B_X(A)\), the transpose
  \(\mathbb D(q_A)\), and a finite-type or completed pair map
  \(\nu_A\colon K_X(A)\to A^!\);
\item a specified binary collision class
  \(r_A^{\mathrm{coll}}\) in its declared target.
\end{enumerate}
A morphism of triples contains maps on every displayed component.
\end{definition}

\begin{construction}[DS comparison on modular Koszul triples]
\label{constr:ds-functor-triples}
\ClaimStatusOpen
Apply \(\mathbf R\mathrm{DS}_f\) to the algebra component and the
comparison map \(\eta_{A,f}\) to the bar component.  Verdier
compatibility then gives the candidate map on \(A^!\), while the
binary projection of the chain map gives the candidate collision
class.  Verifying these maps supplies the DS functor on triples.
\end{construction}

\begin{theorem}[BRST reduction on a full modular datum]
\label{thm:ds-platonic-functor}
\ClaimStatusConditional
Assume DS--bar exchange, BRST concentration, Verdier compatibility,
determinant-line compatibility, and preservation of the declared
collision and modular classes.  Then BRST reduction defines a
morphism of the corresponding full modular data.  Each component is
the map supplied by its named hypothesis.
\end{theorem}

\begin{remark}[Derived reduction precedes invariant extraction]
\label{rem:ds-before-shadows}
The derived BRST complex is the source object.  Scalar, collision,
and modular invariants are extracted after the comparison maps have
been constructed.  This order records the ghost and neutral-field
contributions in their natural chain complex.
\end{remark}

\begin{corollary}[Maurer--Cartan descent]
\label{cor:ds-theta-descent}
\ClaimStatusConditional
Under Theorem~\ref{thm:ds-platonic-functor}, a Maurer--Cartan element
\(\Theta_A\) descends when the comparison map is compatible with the
deformation \(L_\infty\)-algebras:
\[
\Theta_{\mathbf R\mathrm{DS}_f(A)}
=\mathbf R\mathrm{DS}_f(\Theta_A).
\]
\end{corollary}

\begin{remark}[Non-principal scope]
\label{rem:ds-triple-nonprincipal}
The BP family is the first non-principal test.  Its exact OPE and
central data are known; its genus-one curvature and DS--bar map are
the first unresolved entries of the modular triple.
\end{remark}

\subsection{Large-rank character stabilization}
\label{sec:w-infty}
\label{sec:w-completion-kinematics}

The generic universal principal algebra has the vacuum character
\begin{equation}
\label{eq:wn-generic-character-deep}
\chi_N(q)
=\prod_{s=2}^{N}\prod_{m\ge0}(1-q^{s+m})^{-1}
=\prod_{n\ge2}(1-q^n)^{-\min(n-1,N-1)}.
\end{equation}
This is a state-counting theorem.  It becomes a bar statement after
the bar differential has been computed.

\begin{definition}[Stable principal character datum]
\label{def:w-infty}
\ClaimStatusDefinitional
The coefficientwise stable character is
\[
\chi_{\mathrm{st}}(q)
:=\lim_{N\to\infty}\chi_N(q)
=\prod_{n\ge2}(1-q^n)^{-(n-1)}.
\]
This definition introduces a formal character datum.  A vertex
algebra \(\mathcal W_\infty\) additionally requires compatible
algebra maps and a completed category.
\end{definition}

\begin{proposition}[Generic principal character]
\label{prop:wn-character-primitive}
\ClaimStatusProvedHere
Equation~\eqref{eq:wn-generic-character-deep} follows from free strong
generation in weights \(2,\ldots,N\).  Its single-particle series is
\begin{equation}
\label{eq:primitive-series}
P_N(q):=\sum_{s=2}^{N}\frac{q^s}{1-q},
\end{equation}
and the character identity is
\begin{equation}
\label{eq:free-algebra-inversion}
\chi_N(q)=\operatorname{PE}[P_N(q)].
\end{equation}
\end{proposition}

\begin{computation}[Coefficient stabilization]
\label{comp:wn-stabilization-windows}
\ClaimStatusProvedHere
For fixed \(m\), the coefficient \([q^m]\chi_N(q)\) is constant for
\(N\ge m\), since generators of weight \(s>m\) contribute above the
window.  Consequently
\[
[q^m]\chi_{\mathrm{st}}(q)=[q^m]\chi_m(q).
\]
\end{computation}

\begin{proposition}[MacMahon factorization]
\label{prop:winfty-macmahon-deep}
\ClaimStatusProvedHere
Let
\[
M(q)=\prod_{n\ge1}(1-q^n)^{-n},
\qquad
\varphi(q)=\prod_{n\ge1}(1-q^n).
\]
Then
\begin{equation}
\label{eq:bar-window-series}
\chi_{\mathrm{st}}(q)=M(q)\varphi(q).
\end{equation}
\end{proposition}

\begin{proof}
Multiplication adds exponents \(-n+1=-(n-1)\); the \(n=1\)
factor has exponent zero.
\end{proof}

\begin{conjecture}[A completed principal tower]
\label{conj:w-infty-bar}
\ClaimStatusConjectured
Choose compatible principal-algebra maps, a weightwise complete
category, and a pro-conilpotent ordered-bar functor.  The resulting
pro-object should admit a convergent comparison
\[
\bar B_X\!\left(\varprojlim_N\mathcal W_N\right)
\longrightarrow
\varprojlim_N\bar B_X(\mathcal W_N).
\]
The proof requires explicit transition maps and control of the
derived inverse-limit term.
\end{conjecture}

\begin{theorem}[Large-rank factorization/Koszul comparison package]
\label{thm:winfty-factorization-kd}
\ClaimStatusOpen
A large-rank Koszul theorem requires:
\begin{enumerate}[label=\textup{(W\arabic*)}]
\item a defined pro-vertex-algebra and compatible OPE maps;
\item coefficientwise complete bar complexes;
\item vanishing or computation of the relevant
  \(\varprojlim^{1}\)-terms;
\item continuous Verdier duality and cobar convergence.
\end{enumerate}
Under these hypotheses, the desired comparison is the inverse limit
of the finite-stage comparison maps.
\end{theorem}

\begin{theorem}[Stable scalar cyclic line obligation]
\label{thm:winfty-scalar}
\ClaimStatusOpen
The one-dimensionality of a stable scalar cyclic line requires a
defined inverse system of cyclic complexes and an isomorphism theorem
for their degree-two scalar subquotients.  The character stabilization
above supplies the finite-weight state counts; the cyclic differential
and transition maps supply the remaining proof.
\end{theorem}

\begin{remark}[Bar-growth quantities]
\label{prop:virasoro-primitive}
\label{prop:virasoro-entropy}
\label{comp:w3-naive-corrected}
\label{prop:w3-entropy}
\label{prop:wn-entropy-ladder}
\label{prop:winfty-bar-window}
\ClaimStatusOpen
For an ordered-bar complex with finite weight windows, define
\begin{equation}
\label{eq:virasoro-primitive}
\label{eq:w3-naive-primitive}
\label{eq:w3-naive-window}
\label{eq:w3-lambda-coproduct}
\label{eq:w3-corrected-primitive}
\label{eq:w3-corrected-window}
\label{eq:w3-corrected-recursion}
G_A^{\mathrm{bar}}(q)
:=\sum_{m,r}\dim
\operatorname{Prim}H^r\!\left(\bar B_X(A)\right)_m q^m,
\qquad
H_A^{\mathrm{bar}}(q)
:=\sum_{m,r}\dim H^r\!\left(\bar B_X(A)\right)_m q^m.
\end{equation}
These series are determined by the bar differential.  Vacuum
characters determine the chain-space census, while OPE collision
maps and configuration-space boundaries determine cohomology.
Growth radii and entropy are defined after these coefficients have
been computed.
\end{remark}

\begin{remark}[Principal shadow data]
\label{sec:higher-w-shadow-data}
The universal \(W_3\) OPE and the Virasoro Gram matrix give the exact
finite collision packet.  The higher shadow tower is an unbounded
transfer problem in the modular deformation complex.  The transfer
and convergence hypotheses determine its theorem status; the vacuum
character determines the state-counting input.
\end{remark}

\subsection{Corner VOAs and adjacent superconformal lanes}
\label{sec:y-algebra-koszulness}

\begin{theorem}[Chiral Koszulness of \(Y\)-algebras]
\label{thm:y-algebra-koszulness}
\ClaimStatusOpen
For a Gaiotto--Rap\v{c}\'ak corner VOA
\(Y_{N_1,N_2,N_3}[\Psi]\), a chiral Koszul theorem requires:
a proved generic strong-generation presentation, a comparison from
its PBW associated graded to the ordered chiral bar construction,
spectral-sequence convergence, and control of truncation ideals at
special \(\Psi\).  The BRST and triality presentations supply
candidate comparison routes.  Their completion into the displayed
bar theorem remains open.
\end{theorem}

\begin{remark}[Parameter loci]
\label{rem:y-algebra-non-generic-locus}
Generic and truncation parameters are treated in separate families.
At each exceptional parameter, null fields alter the bar differential
and require a quotient comparison.
\end{remark}

\begin{remark}[Triality and duality]
\label{rem:y-algebra-koszul-duality-compat}
Triality is an isomorphism of the appropriately labelled corner
VOAs.  A Koszul-dual interpretation of a parameter reflection
requires the completed bar--Verdier comparison.
\end{remark}

\begin{remark}[Depth classification]
\label{rem:y-algebra-depth-classification-deep-summary}
Shadow depth and Koszulness are independent theorem coordinates:
the first concerns transferred modular operations, while the second
concerns bar cohomology.
\end{remark}

\subsection{The \(N=2\) comparison}
\label{sec:n2-shadow-tower}

\begin{remark}[\(N=2\) superconformal algebra]
\label{thm:n2-kappa}
\label{rem:n2-sca-koszulness}
\label{rem:n2-spectral-flow-invariance}
Chapter~\ref{chap:n2-sca} gives the \(N=2\) OPE, spectral-flow, and
bar-comparison calculations.  The scalar identity used here is
\begin{equation}
\label{eq:n2-complementarity}
K^\kappa_{N=2}
:=\kappa_{N=2}+\kappa_{N=2}^{!}.
\end{equation}
Its evaluation belongs to the companion chapter's stated hypothesis
package.
\end{remark}

\begin{conjecture}[Pixton comparison from a modular MC tower]
\label{conj:pixton-from-shadows}
\ClaimStatusConjectured
A comparison with Pixton relations requires a constructed morphism
from the modular deformation complex to the tautological
Feynman-transform complex, compatibility with the unit and gluing,
and a generation theorem for its image.  Finite-genus calculations
provide evidence after these maps are fixed.
\end{conjecture}

\subsection{Arithmetic recognition kept on its own surface}
\label{sec:walgdeep-supplement}

The Borcherds singular-theta theorem assigns weight
\(\frac12 f(0,0)\) to a multiplicative lift from a weight-zero
Jacobi input with Fourier coefficient \(f(0,0)\)
\cite[Theorem~13.3]{Borcherds1998Grassmannian}.  The input Jacobi
form, its lattice, its polar part, and its character are therefore
load-bearing data.

\begin{remark}[Borcherds and additive Gritsenko weights]
\label{rem:walgdeep-W16-borch-vs-grit}
The multiplicative Borcherds weight is \(f(0,0)/2\).
An additive Gritsenko lift carries the weight of its Jacobi cusp-form
input.  For the standard index-one generator
\(\phi_{0,1}=y+10+y^{-1}+O(q)\), whose double is the K3 elliptic
genus, the coefficient \(f(0,0)=10\) gives Borcherds weight~\(5\).
\end{remark}

\begin{definition}[Schur--Jacobi/Borcherds recognition package]
\label{def:walgdeep-schur-jacobi-borcherds-recognition}
\ClaimStatusDefinitional
A recognition package \(\mathcal R_N\) consists of:
\begin{enumerate}[label=\textup{(R\arabic*)}]
\item an explicitly constructed Weyl-invariant weak lattice-Jacobi
  form \(\phi^{(N)}\);
\item a proof of its modular law, lattice index, and Fourier
  integrality;
\item its complete polar part and the exact coefficient
  \(f^{(N)}(0,0)\);
\item the Borcherds product with its Weyl vector, divisor, character,
  and normalization;
\item a map from the proposed Schur or VOA source to the product's
  graded root spaces, compatible with brackets and completion.
\end{enumerate}
\end{definition}

\begin{theorem}[Weight implied by a recognition package]
\label{thm:walgdeep-gaiotto-siegel-weight}
\ClaimStatusConditional
Assume \(\mathcal R_N\) and the verified coefficient identity
\[
f^{(N)}(0,0)=2(N+3).
\]
Then Borcherds' theorem gives multiplicative-lift weight
\[
\operatorname{wt}\operatorname{Borch}(\phi^{(N)})=N+3.
\]
Every class-\(\mathcal S\), Niemeier, or moonshine interpretation is
carried by the maps and equivariance data in \(\mathcal R_N\).
\end{theorem}

\begin{proof}
Apply Borcherds' weight formula
\(\operatorname{wt}=f^{(N)}(0,0)/2\).
\end{proof}

\begin{remark}[The \(N=3,4\) coefficient obligation]
\label{rem:walgdeep-W17-borcherds-kN-computation}
\ClaimStatusOpen
The values \(f^{(3)}(0,0)\) and \(f^{(4)}(0,0)\) require explicit
Schur/Jacobi inputs and coefficient extraction.  An adjoint
dimension and a rank count supply consistency checks; the Jacobi
form supplies the coefficient.
\end{remark}

\begin{theorem}[Pure \(A\)-type rank arithmetic]
\label{thm:walgdeep-divisor-rule}
\ClaimStatusProvedHere
A root system \(mA_{N-1}\) has rank \(m(N-1)\).  It has rank \(24\)
exactly when
\[
(N-1)\mid24,\qquad m=\frac{24}{N-1}.
\]
For \(2\le N\le25\), the solutions are
\[
N\in\{2,3,4,5,7,9,13,25\}.
\]
\end{theorem}

\begin{proof}
This is the divisibility condition for \(m(N-1)=24\).
\end{proof}

\begin{remark}[Niemeier classification data at Coxeter number \(6\)]
\label{rem:walgdeep-W18-6D4-reanchoring}
\label{thm:walgdeep-N6-reanchoring}
\label{thm:walgdeep-N6-reanchor-A5-4-D4}
\ClaimStatusOpen
The Niemeier root systems \(A_5^4D_4\) and \(D_4^6\) share Coxeter
number \(6\); the first contains \(A_5\) summands.  A class-\(\mathcal
S\) anchor requires an explicit lattice map from the flavour data and
an equivariant Jacobi-form comparison.  The classification fact
provides the candidate list, while the recognition package provides
the anchor theorem.
\end{remark}

\begin{remark}[Higher-\(N\) Niemeier comparisons]
\label{thm:walgdeep-N7-N8-re-anchor}
\label{thm:walgdeep-substitute-anchors}
\label{rem:walgdeep-substitute-anchors}
\label{rem:walgdeep-coxeter-void-regime}
\ClaimStatusOpen
Rank and Coxeter-number tables classify candidate Niemeier root
systems.  A sequence of class-\(\mathcal S\) anchors additionally
requires the family of Jacobi inputs \(\phi^{(N)}\), their lattice
maps, and compatible equivariance.  Definition~\ref{def:walgdeep-schur-jacobi-borcherds-recognition}
records this finite recognition problem at each \(N\).
\end{remark}

\begin{theorem}[Leech lattice fact and comparison obligation]
\label{thm:walgdeep-N24-conway}
\ClaimStatusConditional
The Leech lattice is the unique positive-definite even unimodular
rank-\(24\) lattice with empty root system, and its full automorphism
group is \(\mathrm{Co}_0\).  A comparison with an \(A_{23}\)
class-\(\mathcal S\) Schur sector requires a map from that sector to
the Leech or Conway module together with character and product
compatibility.  Under such a recognition map, Conway equivariance
becomes a property of the source.
\end{theorem}

\begin{theorem}[Extended weight-ladder obligation]
\label{thm:walgdeep-N13-N24-ladder}
\ClaimStatusOpen
For each \(N\), Theorem~\ref{thm:walgdeep-gaiotto-siegel-weight}
produces weight \(N+3\) from the hypothesis
\(f^{(N)}(0,0)=2(N+3)\).  A ladder theorem over a range of \(N\)
therefore consists of explicit recognition packages
\(\{\mathcal R_N\}\) and compatibility maps between them.
\end{theorem}

\subsection{Status ledger}

\begin{center}
\small
\renewcommand{\arraystretch}{1.25}
\begin{tabular}{p{0.30\textwidth}p{0.18\textwidth}p{0.42\textwidth}}
\toprule
Claim & Status & Proof surface \\
\midrule
Principal \(c_N(k)\) and reflected sum
& proved
& Fateev--Lukyanov formula plus \(t\mapsto-t\) \\
Principal \(\kappa=c(H_N-1)\)
& conditional
& \(H_{\mathrm{genus1}}^{W_N}\) \\
BP parity and OPE packet
& proved elsewhere
& FKR, Definition~2.1 and OPE normalization \\
BP standard central sum \(50\)
& proved
& rational identity in \(t=k+3\) \\
BP shifted sum \(196\)
& computed-secondary
& shifted conformal-vector rational function \\
BP \(17/6\)
& diagnostic
& even reciprocal-weight sum \\
BP \(\kappa,\rho,K^\kappa\)
& open
& genus-one BRST curvature \\
BP same-family duality
& open
& DS--bar--Verdier comparison \\
Arbitrary-nilpotent DS/bar commutation
& conjectured
& filtered comparison map and convergence \\
\(W_3\) degree-three bar cohomology
& open
& finite differential-matrix rank computation \\
Large-\(N\) vacuum character
& proved
& coefficientwise free-generation count \\
Large-\(N\) bar/Koszul limit
& open
& pro-system, \(\varprojlim^1\), continuous duality \\
Class-\(\mathcal S\)/Borcherds recognition
& conditional
& package \(\mathcal R_N\) \\
\bottomrule
\end{tabular}
\end{center}
